
Multi-objective code generation requires optimizing multiple quality dimensions simultaneously: correctness (passing tests), security (avoiding vulnerabilities), maintainability (code quality), and efficiency (performance). Current approaches use either weighted multi-objective reinforcement learning with manual weight tuning, or architectural factorization methods that decompose the problem into aspect-specific subspaces with specialized modules.

The architectural factorization approach assumes that quality dimensions can be separated into independent subspaces, justifying complex designs such as gated adapters with aspect-specific routing, LoRA modules with orthogonality constraints, and multi-task architectures with task-specific subspaces. This assumption is borrowed from multi-task learning in vision and natural language processing, where such methods have proven effective.

However, this assumption has not been empirically validated for code. The implicit logic is: if metrics are independent, then code modifications should exhibit directional structure along aspect-aligned axes. This would justify architectural complexity. But independence (uncorrelated metrics) and factorization (aspect-aligned eigenvectors) are distinct concepts. Metrics can be independent yet changes can be multi-directional with no natural coordinate system.

\subsubsection{Research Question}

We address a methodological question: how can we test whether multi-objective code modifications exhibit aspect-dominant structure that would justify architectural factorization? We present a validation framework combining spectral analysis, permutation testing, and cross-validation.

\subsubsection{Contributions}

This work makes four methodological contributions:

\begin{enumerate}
\item \textbf{Validation framework}: We introduce a multi-angle validation approach combining residual covariance analysis, spectral decomposition, permutation testing with 1000 shuffles, directional stability measures, and leave-one-repo-out cross-validation. This framework prevents false positives from noise overfitting.
\end{enumerate}

\begin{enumerate}
\item \textbf{Conceptual distinction}: We clarify the difference between independence and factorization through proof-of-concept demonstration. Uncorrelated metrics (diagonal-dominant covariance) do not automatically yield aspect-aligned eigenvectors. Code modifications can be spherically distributed even when measurements are orthogonal.
\end{enumerate}

\begin{enumerate}
\item \textbf{Pipeline validation}: Through synthetic test data evaluation, we demonstrate that the framework correctly identifies the absence of factorization when it does not exist. The synthetic data showed independence (coupling = 0.072) without directional structure (gap = 1.580, p = 0.955).
\end{enumerate}

\begin{enumerate}
\item \textbf{Alternative hypotheses}: We propose four testable structural hypotheses for future empirical work: hierarchical quality (DAG dependencies), contextual factorization (domain-specific), commit size dependence (scale-dependent), and temporal dynamics (sequential patterns).
\end{enumerate}

\subsubsection{Critical Limitation}

All experimental results presented use synthetic test data generated for pipeline validation, NOT real GitHub commits. The synthetic data was designed to demonstrate that our methodology correctly identifies lack of factorization when none exists. Results have zero scientific validity for claims about real developer behavior. Real data collection is required before conclusions about actual development practices can be drawn. We emphasize this limitation transparently throughout.
